\section*{Bab \ref{ch:g9:natural-risks} --- Hidup Bersama Risiko Alam}

\begin{solution}{exo:g9:natural-risks:1}
(a) Bahaya. (b) \omterm{def:g9:natural-risks:vulnerability}{Kerentanan}. (c) \omterm{def:g9:natural-risks:vulnerability}{Kerentanan}. (d) Bahaya. (e) Keduanya:
aliran lumpur adalah bahayanya; membangun kota di jalurnya membuat kota
itu rentan.
\end{solution}

\begin{solution}{exo:g9:natural-risks:2}
(a) Salah: di tempat yang tidak berpenghuni, tidak ada yang dapat
dicelakai, sehingga \omterm{def:g9:natural-risks:risk}{risikonya} rendah. (b) Benar. (c) Benar. (d) Salah:
peringatan tidak mengubah \omterm{def:g4:volcanoes-earthquakes-first-look:volcano}{gunung apinya}; peringatan menurunkan \omterm{def:g9:natural-risks:vulnerability}{kerentanan}
penduduk, yang dapat pergi tepat waktu.
\end{solution}

\begin{solution}{exo:g9:natural-risks:3}
(a) Ya, beberapa hari sebelumnya. (b) Ya, beberapa jam sampai beberapa
hari sebelumnya. (c) Tidak: \omterm{def:g4:volcanoes-earthquakes-first-look:earthquake}{gempa bumi} tidak dapat diramalkan. (d) Ya:
setelah \omterm{def:g7:earthquakes:tsunami}{tsunami} terjadi, \omterm{def:g7:earthquakes:tsunami}{tsunami} itu memerlukan berjam-jam untuk
menyeberangi \omterm{def:g2:water-on-earth:ocean}{samudra}, dan peringatan merambat lebih cepat.
\end{solution}

\begin{solution}{exo:g9:natural-risks:4}
(a) Guncangan tanah: rentetan gempa-gempa kecil selagi \omterm{def:g7:volcanoes:magma}{magma} naik.
(b) Kemiringan lereng: \omterm{def:g4:volcanoes-earthquakes-first-look:volcano}{gunung api} yang menggembung. (c) Bagaimana tanah
bergeser, beberapa milimeter: penggembungan selagi \omterm{def:g7:volcanoes:magma}{magma} berkumpul.
(d) Gas-gas dari kepulan asap: makin banyak gas yang keluar selagi \omterm{def:g7:volcanoes:magma}{magma}
naik.
\end{solution}

\begin{solution}{exo:g9:natural-risks:5}
(a) $45 \div 2 = \qty{22.5}{km/h}$. (b) $4 \times 2 = \qty{8}{km}$, cukup
jauh untuk mencapai tempat yang lebih tinggi. (c) Tak seorang pun melihat
\omterm{def:g4:volcanoes-earthquakes-first-look:eruption}{letusan} itu, yang tersembunyi oleh \omterm{def:g5:natural-hazards:storm}{badai}, dan tak seorang pun
memperingatkan mereka: tidak ada cara cepat untuk mengirimkan keterangan
ke kota itu, dan tidak ada rencana evakuasi. (d) Misalnya: rencana
evakuasi dengan jalur ke tempat yang lebih tinggi, yang dilatih
sebelumnya; rantai peringatan dari para ilmuwan sampai ke kota; lebih
banyak peralatan di \omterm{def:g4:volcanoes-earthquakes-first-look:volcano}{gunung api} itu.
\end{solution}

\begin{solution}{exo:g9:natural-risks:6}
(a) \omterm{def:g5:water-shapes-the-land:flood}{Banjir} \qty{10}{m} di atas \omterm{def:g1:landscapes-around-us:water}{sungai}. (b) Sekitar \qty{0.47}{km} untuk
\omterm{def:g5:water-shapes-the-land:flood}{banjir} \qty{2}{m}, \qty{0.56}{km} untuk \omterm{def:g5:water-shapes-the-land:flood}{banjir} \qty{5}{m}, dan
\qty{0.67}{km} untuk \omterm{def:g5:water-shapes-the-land:flood}{banjir} \qty{10}{m} (pembacaan dengan selisih sekitar
\qty{0.05}{km} dapat diterima). (c) Hanya \omterm{def:g5:water-shapes-the-land:flood}{banjir} \qty{10}{m}: tanahnya
lebih tinggi daripada 2 dan \qty{5}{m} di atas \omterm{def:g1:landscapes-around-us:water}{sungai}. (d) Hanya
diperlukan sedikit kenaikan air \omterm{def:g1:landscapes-around-us:water}{sungai} untuk mencapainya, dan \omterm{def:g5:water-shapes-the-land:flood}{banjir} kecil
lebih sering datang daripada \omterm{def:g5:water-shapes-the-land:flood}{banjir} besar.
\end{solution}

\begin{solution}{exo:g9:natural-risks:7}
(a) Bahaya tinggi, \omterm{def:g9:natural-risks:vulnerability}{kerentanan} rendah: \omterm{def:g9:natural-risks:risk}{risiko} sedang. (b) Bahaya rendah,
\omterm{def:g9:natural-risks:vulnerability}{kerentanan} tinggi: \omterm{def:g9:natural-risks:risk}{risiko} sedang. (c) Bahaya tinggi, \omterm{def:g9:natural-risks:vulnerability}{kerentanan} tinggi:
\omterm{def:g9:natural-risks:risk}{risiko} tinggi. (d) Bahaya rendah, \omterm{def:g9:natural-risks:vulnerability}{kerentanan} rendah: \omterm{def:g9:natural-risks:risk}{risiko} rendah.
\end{solution}

\begin{solution}{exo:g9:natural-risks:8}
(a) C, lalu B, lalu A. (b) Penguat diagonal membentuk segitiga, yang tidak
dapat miring tanpa meregangkan atau memampatkan sebuah penguat: rangka itu
tetap mempertahankan bentuknya alih-alih miring. (c) Isolator bergeser
mengikuti gerakan tanah yang maju mundur di bawahnya, sehingga jauh lebih
sedikit gerakan yang sampai ke sekolah.
\end{solution}

\begin{solution}{exo:g9:natural-risks:9}
(a) $1 \div 100 = 1~\%$ per tahun. (b) Peluangnya sama setiap tahun, apa
pun yang terjadi pada tahun sebelumnya: \omterm{def:g5:water-shapes-the-land:flood}{banjir} itu dapat datang lagi
tahun depan. (c) $300 \div 100 = 3$ \omterm{def:g5:water-shapes-the-land:flood}{banjir}, rata-rata.
\end{solution}

\begin{solution}{exo:g9:natural-risks:10}
(a) $3600 \div 15 = 240$ kali pembacaan. (b) $240 \times 24 \times 365 =
\num{2102400} \approx 2.1 \times 10^6$ kali pembacaan. (c) $2400 \div 800 =
\qty{3}{h}$. (d) Dalam tiga jam itu, pusat-pusat peringatan dapat
memperingatkan pantai dan penduduk dapat pindah ke tempat yang lebih
tinggi.
\end{solution}

\begin{solution}{exo:g9:natural-risks:11}
(a) $4 \times 3 \times 3.8 = \qty{45.6}{L}$. (b) $45.6 \div 1.5 = 30.4$,
jadi 31 botol. (c) Misalnya: makanan yang tahan lama, untuk dimakan ketika
toko-toko tutup; senter, jika listrik padam; radio yang bekerja dengan
baterai atau diputar dengan tangan, untuk mendengar peringatan; kotak P3K,
untuk merawat luka ringan; selimut hangat. (d) Agar tas itu dapat langsung
dibawa jika keluarga itu harus pergi dengan cepat.
\end{solution}

\begin{solution}{exo:g9:natural-risks:12}
Misalnya: (1) menggambar \omterm{def:g9:natural-risks:hazard-map}{peta bahaya} aliran lumpur, \omterm{def:g1:landscapes-around-us:landform}{lembah} demi \omterm{def:g1:landscapes-around-us:landform}{lembah}
(sisi bahaya); (2) memasang \omterm{def:g7:earthquakes:seismometer}{seismometer}, tiltmeter, penerima GPS, dan alat
ukur gas di \omterm{def:g4:volcanoes-earthquakes-first-look:volcano}{gunung api} itu (sisi bahaya); (3) menyiapkan rantai peringatan
yang cepat dari observatorium sampai ke kota, dengan sirene (sisi
\omterm{def:g9:natural-risks:vulnerability}{kerentanan}); (4) merencanakan jalur ke tempat yang lebih tinggi,
memasang rambunya, dan melatihnya dengan latihan rutin (sisi \omterm{def:g9:natural-risks:vulnerability}{kerentanan});
(5) tidak membangun rumah atau sekolah baru di jalur aliran lumpur (sisi
\omterm{def:g9:natural-risks:vulnerability}{kerentanan}).
\end{solution}

\begin{solution}{pb:g9:natural-risks:1}
\textbf{1.} \omterm{def:g1:landscapes-around-us:water}{Sungai} berada pada \qty{143}{m} di samping C, \qty{137}{m} di
samping A, dan \qty{126}{m} di samping B; pada \omterm{def:g2:drawing-to-map:map}{peta}, C berada di barat dan
B di timur, sehingga \omterm{def:g1:landscapes-around-us:water}{sungai} itu mengalir ke timur.

\textbf{2.} A: $137 - 137 = \qty{0}{m}$; B: $276 - 126 = \qty{150}{m}$;
C: $236 - 143 = \qty{93}{m}$.

\textbf{3.} Hanya A, yang \omterm{def:g4:world-physical-map:altitude}{ketinggiannya} kurang dari \qty{5}{m} di atas \omterm{def:g1:landscapes-around-us:water}{sungai}.

\textbf{4.} Sekitar \qty{0.57}{km} (antara 0.5 dan \qty{0.65}{km}).

\textbf{5.} Kenaikan sepanjang \qty{100}{m} adalah separuh kenaikan
sepanjang \qty{200}{m}: A $1.4 \div 2 = 0.7~\%$; B $83 \div 2 = 41.5~\%$;
C $20 \div 2 = 10~\%$.

\textbf{6.} B: lerengnya sangat curam, dan \omterm{def:g3:stones-pebbles-sand:rock}{batuan} serta tanah yang lepas
dapat longsor menuruni lereng yang curam, terutama setelah \omterm{def:g4:clouds-rain-snow:precipitation}{hujan} lebat.

\textbf{7.} A: sekitar \qty{1}{km}; B: hampir tepat di atas \omterm{def:g7:earthquakes:fault}{sesar} (kurang
dari \qty{0.05}{km}); C: sekitar \qty{1.3}{km}.

\textbf{8.} Jika \omterm{def:g7:earthquakes:fault}{sesar} itu bergerak, tanah patah dan bergeser di
sepanjangnya, yang dapat merobek sebuah bangunan; selain itu, gempa pada
\omterm{def:g7:earthquakes:fault}{sesar} itu akan mengguncang tanah di sana dengan sangat keras.

\textbf{9.}

\begin{center}
{\small
\begin{tabular}{|c|c|c|c|}
\hline
lokasi & di zona \omterm{def:g5:water-shapes-the-land:flood}{banjir}? & lereng curam? & di atas \omterm{def:g7:earthquakes:fault}{sesar}? \\
\hline
A & ya & tidak & tidak \\
B & tidak & ya & ya \\
C & tidak & tidak & tidak \\
\hline
\end{tabular}}
\end{center}

\textbf{10.} Lokasi C: hanya lokasi ini yang mendapat tiga kali
``tidak'': \qty{93}{m} di atas \omterm{def:g1:landscapes-around-us:water}{sungai}, lereng yang landai sebesar
\qty{10}{\%}, dan lebih dari \qty{1}{km} dari \omterm{def:g7:earthquakes:fault}{sesar}.

\textbf{11.} Misalnya: penguat diagonal dan dinding geser pada rangkanya,
sambungan yang kuat, atau isolator di bawah bangunan; para murid melatih
tindakan saat gempa (berlutut, berlindung, berpegangan).

\textbf{12.} \textbf{17 rumah} dari 28 berdiri di dalam zona \omterm{def:g5:water-shapes-the-land:flood}{banjir}:
$17 \div 28 \approx 0.61$, kira-kira 3 dari setiap 5 rumah.
\end{solution}
